%% Guided notes, AP Physics 1 Unit 2: Force and Translational Dynamics (topics 2.1 to 2.9). Use g = 10 m/s^2.
\newcommand{\gr}[2]{\begin{center}\begin{tikzpicture}\begin{axis}[gridaxis,#1]#2\end{axis}\end{tikzpicture}\end{center}}
\NewEnviron{nrow}[1]{\begin{tcolorbox}[enhanced,sidebyside,sidebyside align=top seam,lefthand width=2.7cm,sidebyside gap=12pt,colback=white,colframe=apnavy!55,boxrule=0.5pt,arc=0pt,left=5pt,right=6pt,top=5pt,bottom=5pt,before skip=0pt,after skip=-0.5pt,segmentation style={solid,apnavy!35,line width=0.5pt},fontupper=\sffamily\bfseries\color{apnavy}\raggedright]#1\tcblower\BODY\end{tcolorbox}}
\newcommand{\ms}{\,\mathrm{m/s}}
\newcommand{\mss}{\,\mathrm{m/s^2}}
\newcommand{\m}{\,\mathrm{m}}
\newcommand{\s}{\,\mathrm{s}}
\newcommand{\N}{\,\mathrm{N}}
\newcommand{\kg}{\,\mathrm{kg}}
% free-body diagram helpers: \fdot at the origin, \farrow{angle}{length}{label}
\newcommand{\fdot}{\fill[apnavy](0,0) circle(2.2pt);}
\newcommand{\farrow}[3]{\draw[->,very thick,apblue](0,0)--(#1:#2) node[pos=1.18,font=\small,text=apnavy]{#3};}

%%DOC 2.1 | Systems and Center of Mass
\begin{center}\textsf{\small\color{apgrey}Objective: choose a system, tell internal from external forces, and locate the center of mass.}\end{center}
\begin{nrow}{Systems}
A \fillin{system} is the object or group of objects we choose to study. Everything else is the \fillin{environment}.\par
Forces between objects inside the system are \fillin{internal}; forces from outside are \fillin{external}.\par
Only \fillin{external} forces can change the motion of the system as a whole.
\end{nrow}
\begin{nrow}{Center of mass}
The center of mass is the point that moves as if all the mass were there and all the external forces acted there.\par\smallskip
In one dimension: $x_\text{cm}=$ \fillin{$\dfrac{m_1x_1+m_2x_2+\cdots}{m_1+m_2+\cdots}$}\par\smallskip
For a uniform symmetric object (rod, ball, box), the center of mass is at its \fillin{geometric center}.
\end{nrow}
\begin{nrow}{Example}
\begin{center}\begin{tikzpicture}[x=1cm]
\draw[thick](0,0)--(4,0);\foreach \i in {0,1,2,3,4}{\draw(\i,0.08)--(\i,-0.08) node[below,font=\scriptsize]{\i\,m};}
\fill[apblue](0,0.25) circle(5pt) node[above=5pt,font=\small]{$2\kg$};\fill[apred](4,0.35) circle(8pt) node[above=7pt,font=\small]{$6\kg$};
\ifkey\draw[->,thick,apgreen](3,-0.9)--(3,-0.3);\node[apgreen,font=\small] at (3,-1.1){cm};\fi
\end{tikzpicture}\end{center}
$x_\text{cm}=$ \fillin{$\dfrac{2(0)+6(4)}{2+6}=3\m$}: the center of mass is closer to the \fillin{heavier} mass.
\end{nrow}
\begin{nrow}{Motion of the center of mass}
A firework explodes in the air. The pieces fly apart because of \fillin{internal} forces, so the center of mass keeps following the \fillin{same parabola} (no air resistance).\par
Two skaters at rest push each other apart: their center of mass \fillin{stays at rest}.
\end{nrow}
\begin{nrow}{Summary}
\ifkey{\color{apnavy}Choose a system; only external forces change its overall motion. The center of mass is the mass-weighted average position, $x_\text{cm}=\sum m_ix_i/\sum m_i$. When internal structure does not matter, a system can be modeled as a single object at its center of mass.}\else\vspace{2.2cm}\fi
\end{nrow}

%%DOC 2.2 | Forces and Free-Body Diagrams
\begin{center}\textsf{\small\color{apgrey}Objective: name the forces on an object and draw a correct free-body diagram.}\end{center}
\begin{nrow}{Forces}
A force is a \fillin{push or pull} between two objects. It is a \fillin{vector}, measured in \fillin{newtons} (N): $1\N=1\kg\cdot\mathrm{m/s^2}$.\par\smallskip
\begin{tabular}{@{}p{0.47\linewidth}p{0.47\linewidth}@{}}
\textbf{Contact forces} & \textbf{Field (non-contact) forces}\\
normal $F_N$, \fillin{friction} $f$, \fillin{tension} $F_T$, spring $F_s$, air resistance & \fillin{gravitational} force $F_g=mg$, electric, magnetic
\end{tabular}
\end{nrow}
\begin{nrow}{Main forces}
\fillin{Weight} $F_g=mg$: always straight down. \quad A $3\kg$ mass: $F_g=$ \fillin{$30\N$}.\par
\fillin{Normal} force: perpendicular to the surface, pushing \fillin{away from} it.\par
\fillin{Tension}: along the rope, \fillin{pulling} on the object. \quad \fillin{Friction}: parallel to the surface, opposing sliding.
\end{nrow}
\begin{nrow}{Rules for a free-body diagram}
1) Draw the object as a \fillin{dot}. \quad 2) Draw only forces exerted \fillin{on} the object, starting at the dot.\par
3) Arrow lengths show the relative \fillin{sizes}. \quad 4) Label each force with its \fillin{type} and source.\par
5) Never draw \fillin{velocity} or \fillin{acceleration} arrows, and no ``force of motion''.
\end{nrow}
\begin{nrow}{Examples}
\begin{tabular}{@{}p{0.3\linewidth}p{0.3\linewidth}p{0.33\linewidth}@{}}
Book at rest on a table & Lamp hanging from a cord & Box pulled right on a rough floor\\
\begin{tikzpicture}[scale=0.9]\ifkey\farrow{90}{0.9}{$F_N$}\farrow{270}{0.9}{$F_g$}\fi\fdot\end{tikzpicture} &
\begin{tikzpicture}[scale=0.9]\ifkey\farrow{90}{0.9}{$F_T$}\farrow{270}{0.9}{$F_g$}\fi\fdot\end{tikzpicture} &
\begin{tikzpicture}[scale=0.9]\ifkey\farrow{90}{0.8}{$F_N$}\farrow{270}{0.8}{$F_g$}\farrow{0}{1.1}{$F_T$}\farrow{180}{0.6}{$f$}\fi\fdot\end{tikzpicture}
\end{tabular}\par\smallskip
A ball thrown upward, after it leaves the hand (no air resistance): the only force is \fillin{gravity}.
\end{nrow}
\begin{nrow}{Summary}
\ifkey{\color{apnavy}A force is an interaction between two objects. A free-body diagram shows only the forces on one object, as arrows from a dot, labeled and roughly to scale: weight down, normal perpendicular to the surface, tension along the rope, friction along the surface.}\else\vspace{2.2cm}\fi
\end{nrow}

%%DOC 2.3 | Newton's Third Law
\begin{center}\textsf{\small\color{apgrey}Objective: identify third-law force pairs and explain why they do not cancel.}\end{center}
\begin{nrow}{The law}
If object A exerts a force on object B, then B exerts a force on A that is \fillin{equal in magnitude} and \fillin{opposite in direction}:\quad $\vec F_{A\text{ on }B}=$ \fillin{$-\vec F_{B\text{ on }A}$}\par\smallskip
The two forces are the same \fillin{type} and act on \fillin{different} objects, so they \fillin{never cancel} each other.
\end{nrow}
\begin{nrow}{Examples}
Earth pulls an apple down with $1\N$. The apple pulls Earth \fillin{up} with \fillin{$1\N$}.\par
A swimmer pushes the water \fillin{backward}; the water pushes the swimmer \fillin{forward}.\par
A rocket pushes gas \fillin{down}; the gas pushes the rocket \fillin{up}.
\end{nrow}
\begin{nrow}{Equal forces,\\ different effects}
A truck hits a small car. The force on the car is \fillin{equal} to the force on the truck.\par
The car has the \fillin{greater} acceleration, because $a=F/m$ and its mass is \fillin{smaller}.
\end{nrow}
\begin{nrow}{Common mistake}
A book rests on a table. The normal force and the weight are equal, but they are \fillin{not} a third-law pair: they act on the \fillin{same} object and are different types.\par
The partner of the weight (Earth on book) is \fillin{the book pulling Earth up}. The partner of the normal force (table on book) is \fillin{the book pushing the table down}.
\end{nrow}
\begin{nrow}{Summary}
\ifkey{\color{apnavy}Forces come in pairs: equal size, opposite direction, same type, on two different objects. Pair forces never cancel; equal forces can still give different accelerations to objects of different mass.}\else\vspace{2.2cm}\fi
\end{nrow}

%%DOC 2.4 | Newton's First Law
\begin{center}\textsf{\small\color{apgrey}Objective: use the first law and the condition for translational equilibrium.}\end{center}
\begin{nrow}{The law}
If the net force on an object is \fillin{zero}, its velocity stays \fillin{constant}: at rest, it stays at rest; moving, it keeps moving in a straight line at \fillin{constant speed}.\par
\fillin{Inertia}: the tendency of an object to resist changes in its velocity. It is measured by its \fillin{mass}.
\end{nrow}
\begin{nrow}{Equilibrium}
Translational equilibrium: $\sum\vec F=$ \fillin{$0$}, so $\sum F_x=0$ and $\sum F_y=0$.\par
An object in equilibrium can be at rest \textbf{or} moving at \fillin{constant velocity}.
\end{nrow}
\begin{nrow}{Example 1:\\ an elevator}
A $500\kg$ elevator moves up at a constant speed. The cable tension is \fillin{$5000\N$}: the same as the weight, because the acceleration is \fillin{zero}.
\end{nrow}
\begin{nrow}{Example 2:\\ a hanging sign}
\begin{minipage}[t]{0.58\linewidth}
A $100\N$ sign hangs from two ropes, each at $30^\circ$ above the horizontal.\par
Horizontal: the two $T\cos30^\circ$ \fillin{cancel}.\par
Vertical: $2T\sin30^\circ=100\N\Rightarrow T=$ \fillin{$100\N$}.
\end{minipage}\hfill
\begin{minipage}[t]{0.36\linewidth}\vspace{0pt}
\begin{tikzpicture}[scale=0.8]
\draw[thick](-2,1.15)--(0,0)--(2,1.15);\draw[apgrey](-2.3,1.15)--(2.3,1.15);
\fill[apblue!20,draw=apnavy](-0.5,-0.7) rectangle (0.5,0);\node[font=\scriptsize] at (0,-0.35){$100\N$};
\node[font=\scriptsize] at (-1.3,1.0){$30^\circ$};\node[font=\scriptsize] at (1.3,1.0){$30^\circ$};
\end{tikzpicture}
\end{minipage}
\end{nrow}
\begin{nrow}{Check}
A passenger in a bus lurches forward when the bus brakes, because \fillin{his body tends to keep moving} while the bus slows down. No forward force pushes him.
\end{nrow}
\begin{nrow}{Summary}
\ifkey{\color{apnavy}Zero net force means constant velocity (possibly zero). In equilibrium, the forces balance in each direction. Mass measures inertia.}\else\vspace{2.2cm}\fi
\end{nrow}

%%DOC 2.5 | Newton's Second Law
\begin{center}\textsf{\small\color{apgrey}Objective: relate the net force on a system to its acceleration, including connected objects.}\end{center}
\begin{nrow}{The law}
$\vec a=$ \fillin{$\dfrac{\sum\vec F}{m}$}: the acceleration is in the direction of the \fillin{net force}, proportional to it, and inversely proportional to the \fillin{mass}.\par
Strategy: 1) \fillin{free-body diagram} \quad 2) choose an axis along the \fillin{acceleration} \quad 3) write $\sum F=ma$ on each axis.
\end{nrow}
\begin{nrow}{Example 1}
A $4\kg$ box is pushed with $30\N$; friction is $10\N$.\quad $a=$ \fillin{$\dfrac{30-10}{4}=5\mss$}
\end{nrow}
\begin{nrow}{Example 2:\\ apparent weight}
A $60\kg$ person stands on a scale in an elevator accelerating \textbf{up} at $2\mss$.\par
$F_N-mg=ma\Rightarrow F_N=$ \fillin{$60(10+2)=720\N$}: the scale reads \fillin{more} than her weight.\par
Accelerating down at $2\mss$: $F_N=$ \fillin{$60(10-2)=480\N$}. \quad In free fall: $F_N=$ \fillin{$0$}.
\end{nrow}
\begin{nrow}{Example 3:\\ connected objects}
\begin{minipage}[t]{0.6\linewidth}
A $3\kg$ block on a frictionless table is tied to a hanging $2\kg$ block.\par
Whole system: the only unbalanced external force is the weight of the hanging block, \fillin{$20\N$}.\par
$a=$ \fillin{$\dfrac{20}{3+2}=4\mss$}\par
Tension (on the $3\kg$ block): $F_T=$ \fillin{$3\times4=12\N$}
\end{minipage}\hfill
\begin{minipage}[t]{0.34\linewidth}\vspace{0pt}
\begin{tikzpicture}[scale=0.7]
\fill[apgrey!25](0,0) rectangle (3.2,-0.3);\fill[apblue!20,draw=apnavy](0.4,0) rectangle (1.4,0.8);\node[font=\scriptsize] at (0.9,0.4){$3\kg$};
\draw[thick](1.4,0.4)--(3.2,0.4);\draw[thick](3.35,0.25) circle(0.2);\draw[thick](3.55,0.25)--(3.55,-1.2);
\fill[apred!20,draw=apnavy](3.15,-1.2) rectangle (3.95,-1.9);\node[font=\scriptsize] at (3.55,-1.55){$2\kg$};
\end{tikzpicture}
\end{minipage}
\end{nrow}
\begin{nrow}{Summary}
\ifkey{\color{apnavy}$\sum\vec F=m\vec a$. Draw a free-body diagram, pick an axis along the acceleration, and add the forces. For connected objects, treat the whole system first (internal tensions cancel), then one object to find the tension.}\else\vspace{2.2cm}\fi
\end{nrow}

%%DOC 2.6 | Gravitational Force
\begin{center}\textsf{\small\color{apgrey}Objective: use Newton's law of gravitation and the gravitational field.}\end{center}
\begin{nrow}{Universal gravitation}
$F_g=$ \fillin{$G\dfrac{m_1m_2}{r^2}$}, with $G=6.67\times10^{-11}\,\mathrm{N\,m^2/kg^2}$ and $r$ the distance between the \fillin{centers}.\par
\begin{tabular}{@{}p{0.45\linewidth}p{0.45\linewidth}@{}}
$r$ doubles: $F_g$ is \fillin{divided by 4} & one mass triples: $F_g$ is \fillin{multiplied by 3}
\end{tabular}
\end{nrow}
\begin{nrow}{Gravitational field}
$g=\dfrac{F_g}{m}=$ \fillin{$G\dfrac{M}{r^2}$}, in \fillin{N/kg} (the same as m/s$^2$). Near Earth's surface, $g\approx$ \fillin{$10\,\mathrm{N/kg}$}.\par
At a height of one Earth radius above the surface ($r=2R_E$): $g=$ \fillin{$10/4=2.5\,\mathrm{N/kg}$}.
\end{nrow}
\begin{nrow}{Mass vs.\ weight}
Mass (kg) is the \fillin{same} everywhere. Weight $F_g=mg$ depends on \fillin{where} you are.\par
A $60\kg$ astronaut: on Earth $F_g=$ \fillin{$600\N$}; on the Moon ($g=1.6\,\mathrm{N/kg}$) $F_g=$ \fillin{$96\N$}.
\end{nrow}
\begin{nrow}{Weightlessness}
Astronauts in an orbiting station feel weightless because they and the station are \fillin{in free fall} together, not because there is no gravity: at that height $g$ is still about \fillin{$9\,\mathrm{N/kg}$}.
\end{nrow}
\begin{nrow}{Summary}
\ifkey{\color{apnavy}$F_g=Gm_1m_2/r^2$: an inverse-square law. The field $g=GM/r^2$ gives the weight $F_g=mg$. Mass does not change from place to place; weight does. Orbiting objects are in free fall.}\else\vspace{2.2cm}\fi
\end{nrow}

%%DOC 2.7 | Kinetic and Static Friction
\begin{center}\textsf{\small\color{apgrey}Objective: model static and kinetic friction and use them with Newton's laws.}\end{center}
\begin{nrow}{Two kinds}
\textbf{Kinetic} friction (sliding): $f_k=$ \fillin{$\mu_kF_N$}, opposite to the \fillin{sliding} motion.\par
\textbf{Static} friction (no sliding): $f_s\le$ \fillin{$\mu_sF_N$}. It takes \fillin{whatever value} is needed to stop sliding, up to that maximum.\par
Usually $\mu_s$ \fillin{$\ge$} $\mu_k$: it is harder to \fillin{start} an object sliding than to keep it sliding. In this model, friction does \fillin{not} depend on the contact area.
\end{nrow}
\begin{nrow}{Example:\\ a $10\kg$ box}
$\mu_s=0.5$, $\mu_k=0.3$, on a level floor: $F_N=$ \fillin{$100\N$}.\par
Push $20\N$ (box at rest): static friction $=$ \fillin{$20\N$}, the box \fillin{does not move}.\par
To start it moving you need more than \fillin{$50\N$}. Once sliding, friction $=$ \fillin{$30\N$}.\par
Push $40\N$ while it slides: $a=$ \fillin{$\dfrac{40-30}{10}=1\mss$}
\end{nrow}
\begin{nrow}{On an incline}
\begin{minipage}[t]{0.6\linewidth}
A block rests on an incline at angle $\theta$:\par
$F_N=$ \fillin{$mg\cos\theta$} \qquad $f_s=$ \fillin{$mg\sin\theta$}, pointing \fillin{up} the incline.\par
$2\kg$ block, $\theta=30^\circ$: $F_N=$ \fillin{$17.3\N$}, $f_s=$ \fillin{$10\N$}.
\end{minipage}\hfill
\begin{minipage}[t]{0.34\linewidth}\vspace{0pt}
\begin{tikzpicture}[scale=0.75]
\draw[thick,fill=apgrey!15](0,0)--(3.5,0)--(3.5,2.02)--cycle;\node[font=\scriptsize] at (0.75,0.18){$\theta$};
\begin{scope}[shift={(1.9,1.1)},rotate=30]\fill[apblue!20,draw=apnavy](-0.35,0) rectangle (0.35,0.5);\end{scope}
\end{tikzpicture}
\end{minipage}
\end{nrow}
\begin{nrow}{Summary}
\ifkey{\color{apnavy}Kinetic friction $f_k=\mu_kF_N$ opposes sliding. Static friction matches the applied force up to $\mu_sF_N$. Find $F_N$ first (it is not always $mg$), then the friction.}\else\vspace{2.2cm}\fi
\end{nrow}

%%DOC 2.8 | Spring Forces
\begin{center}\textsf{\small\color{apgrey}Objective: use Hooke's law for ideal springs.}\end{center}
\begin{nrow}{Hooke's law}
$\vec F_s=$ \fillin{$-k\,\Delta\vec x$}: the force is proportional to the stretch or compression $\Delta x$ and points back toward the \fillin{equilibrium} (natural) length: a \fillin{restoring} force.\par
$k$ is the \fillin{spring constant}, in \fillin{N/m}: a stiffer spring has a \fillin{larger} $k$.
\end{nrow}
\begin{nrow}{Examples}
$k=200\,\mathrm{N/m}$, stretched $0.10\m$: $F_s=$ \fillin{$20\N$}. \quad Stretched twice as far: \fillin{$40\N$}.\par
A $2\kg$ mass hangs at rest from this spring: $k\Delta x=mg\Rightarrow\Delta x=$ \fillin{$20/200=0.10\m$}.
\end{nrow}
\begin{nrow}{Force vs.\ stretch graph}
\gr{xmin=-0.01,xmax=0.21,ymin=-1,ymax=45,width=6.5cm,height=4.4cm,xlabel={$\Delta x$ (m)},ylabel={$F$ (N)},xtick={0,0.1,0.2},ytick={0,20,40},xticklabel style={/pgf/number format/fixed}}{\addplot[apblue,very thick,domain=0:0.2]{200*x};}
The graph is a \fillin{straight line through the origin}. Its slope is \fillin{$k$}: here \fillin{$200\,\mathrm{N/m}$}.
\end{nrow}
\begin{nrow}{Ideal spring}
An ideal spring has \fillin{no mass}, so it pulls with the \fillin{same} force on the objects at both of its ends.
\end{nrow}
\begin{nrow}{Summary}
\ifkey{\color{apnavy}$F_s=k|\Delta x|$, directed back toward the natural length. $k$ is the slope of the force vs.\ stretch graph. For a hanging mass at rest, $k\Delta x=mg$.}\else\vspace{2.2cm}\fi
\end{nrow}

%%DOC 2.9 | Circular Motion
\begin{center}\textsf{\small\color{apgrey}Objective: find the centripetal acceleration and the net force for an object moving in a circle.}\end{center}
\begin{nrow}{Uniform circular motion}
\begin{minipage}[t]{0.62\linewidth}
The speed is \fillin{constant} but the direction of the velocity \fillin{changes}, so the object \fillin{accelerates}.\par
Centripetal acceleration: $a_c=$ \fillin{$\dfrac{v^2}{r}$}, pointing toward the \fillin{center}.\par
Speed and period: $v=$ \fillin{$\dfrac{2\pi r}{T}$}.
\end{minipage}\hfill
\begin{minipage}[t]{0.32\linewidth}\vspace{0pt}
\begin{tikzpicture}[scale=0.7]
\draw[apgrey,thick](0,0) circle(1.4);\fill(0,0) circle(1.5pt);\fill[apblue](1.4,0) circle(3pt);
\draw[->,very thick,apblue](1.4,0)--(1.4,1.2) node[right,font=\scriptsize]{$\vec v$};
\draw[->,very thick,apred](1.4,0)--(0.4,0) node[above,font=\scriptsize]{$\vec a_c$};
\end{tikzpicture}
\end{minipage}
\end{nrow}
\begin{nrow}{Net force}
$\sum F_\text{toward center}=$ \fillin{$m\dfrac{v^2}{r}$}. It is \fillin{not} a new force: it is provided by tension, gravity, friction or a normal force.\par
If the string breaks, the object moves along the \fillin{tangent}, in a straight line.
\end{nrow}
\begin{nrow}{Examples}
A $0.5\kg$ ball on a string moves in a horizontal circle, $r=1\m$, $v=4\ms$ (ignore gravity): $a_c=$ \fillin{$16\mss$}, $F_T=$ \fillin{$8\N$}.\par
A car on a flat curve: the centripetal force is \fillin{static friction}. $\mu_smg=mv^2/r\Rightarrow v_\text{max}=$ \fillin{$\sqrt{\mu_sgr}$}. With $\mu_s=0.5$, $r=50\m$: \fillin{$\approx15.8\ms$}.\par
Top of a hill of radius $r$: $mg-F_N=mv^2/r$, so the car feels \fillin{lighter}.
\end{nrow}
\begin{nrow}{Summary}
\ifkey{\color{apnavy}In uniform circular motion $a_c=v^2/r$ toward the center, and the net force toward the center is $mv^2/r$, supplied by real forces. Doubling the speed multiplies the required force by 4.}\else\vspace{2.2cm}\fi
\end{nrow}
